Elemento obrigatório (Guia UECE 2026, 2.2.2.1.8, com base na NBR 6028). O resumo é a apresentação concisa dos pontos relevantes do trabalho, dando uma visão rápida e clara do conteúdo e das conclusões. Ele é a última coisa que você escreve, mas a primeira que a banca lê -- capriche.

Ao redigi-lo, siga estas regras:

\begin{alineas}
	\item deve ser informativo: apresente finalidade, metodologia, resultados e conclusões do trabalho;
	\item componha-o com uma sequência de frases concisas e afirmativas, e não com enumeração de tópicos;
	\item use o verbo na voz ativa e na 3ª pessoa do singular;
	\item deve conter de 150 a 500 palavras, em um único parágrafo, sem citações, sem fórmulas e sem abreviações não explicadas;
	\item a primeira frase precisa ser significativa e expressar o tema principal do trabalho -- evite abrir com ``O presente trabalho...'' ou ``O autor descreve...'';
	\item evite frases negativas, símbolos e fórmulas que não sejam de uso corrente, comentário pessoal, crítica ou julgamento de valor;
	\item finalize com as palavras-chave, antecedidas da expressão ``Palavras-chave:'', separadas por ponto e vírgula e finalizadas por ponto, cada uma iniciada por letra minúscula (exceto nomes próprios e científicos).
\end{alineas}

Substitua este parágrafo pelo texto do seu resumo, e a linha abaixo pelas suas próprias palavras-chave (de 3 a 5 termos costuma ser suficiente).

% Separe as palavras-chave por ponto e vírgula ';' e finalizadas por ponto.
\palavraschave{primeira palavra-chave; segunda palavra-chave; terceira palavra-chave.}
